\section{符号说明}

\begin{table}[H]\centering\small
\caption{主要符号、单位与含义}\label{tab:symbols}
\begin{tabular}{clc}\toprule
符号 & 含义 & 单位\\\midrule
$i,j$ & 调度中心或服务区节点 & ---\\
$g$ & 运输机型A、B或C & ---\\
$b$ & 不可拆货箱 & ---\\
$q,Q_g$ & 航段有效载荷、机型额定载荷 & kg\\
$d_{ij},H_{ij}$ & 水平航段长度、巡航海拔 & m\\
$z_i$ & 节点作业海拔 & m\\
$E_g^{\rm use},E_{ij}^g(q)$ & 可用电量、航段能耗 & kWh\\
$\rho_g$ & 返航安全余量比例 & ---\\
$s_p,C_p$ & 架次开始和返航完成时刻 & s\\
$D_b,d_b,f_b$ & 实际交付、期望送达、首批截止时刻 & s\\
$h_b,T_b$ & 适用的硬截止、相对期望时刻的迟到量 & s\\
$\pi_b,\bar D_\pi,T_\pi$ & 优先系数、加权平均交付、加权迟到 & ---、s、加权s\\
$A_{ab}(t)$ & 端点$a,b$的双向链路可用性 & 0或1\\
$R_k$ & 中继悬停位置与海拔 & m及经纬度\\
\bottomrule\end{tabular}
\end{table}

\begin{table}[H]\centering\small
\caption{各问实际使用的决策对象}\label{tab:decisions}
\begin{tabular}{lll}\toprule
问题 & 决策对象 & 含义\\\midrule
一 & $x_a$ & 候选组批是否入选，取0或1\\
二 & $B_p,g_p,s_p$ & 架次箱集、机型与任务开始时刻\\
三 & 中继站点、服务窗口 & 与运输架次共同决定连续通信覆盖\\
四 & 服务区所属组 & 服从不可拆运输任务组件约束\\\bottomrule
\end{tabular}\end{table}
